%%═══════════════════════════════════════════════════════════════════
%% Lecture 08 — Multigrid Method and Introduction to PDEs
%% 77315 — Computational Physics
%% Date: 10/05/2026
%%═══════════════════════════════════════════════════════════════════

\section{Lecture 08 — Multigrid Method and Introduction to
  Partial Differential Equations}\label{sec:lec08}

\begin{center}
\begin{tabular}{ll}
  \textbf{Date:}\quad 10/05/2026 & \\
\end{tabular}
\end{center}
\hrule\vspace{1.5em}

%%──────────────────────────────────────────────────────────────────
\subsection*{Overview}
Lecture~\ref{sec:lec07} introduced the relaxation method for boundary-value
problems; the Gauss–Seidel iteration was identified as the workhorse.
This lecture begins by explaining \emph{why} Gauss–Seidel converges
slowly on large grids and how the \emph{multigrid} algorithm repairs
the deficiency by combining solutions on grids of different spacings.
The second half of the lecture opens a new chapter: \emph{partial
differential equations (PDEs)}. We classify PDEs into elliptic,
parabolic, and hyperbolic types, discuss boundary and initial
conditions, discretise the two-dimensional Poisson equation, and
study the Jacobi, Gauss–Seidel, and successive over-relaxation (SOR)
iterative solvers.

%%──────────────────────────────────────────────────────────────────
\subsection{Motivation: Why Gauss–Seidel Converges Slowly}
\label{subsec:lec08:gs-slow}

\subsubsection*{Physical Motivation}
The Gauss–Seidel iteration at every grid point replaces $u_{i,j}$ by
the average of its four neighbours. This operation is a local smoother:
it rapidly damps oscillations whose wavelength is comparable to the
grid spacing $h$, but it barely touches long-wavelength errors, whose
smoothing rate is proportional to $h^2$.  For an $N\times N$ grid with
$h=1/N$, reaching a fixed accuracy requires $\bigO{N^2}$ iterations,
each costing $\bigO{N^2}$ work, giving a total cost of
$\bigO{N^4}$ — prohibitive even for moderate $N$.

The key observation is that, on a coarser grid with spacing $H=2h$,
the same long-wavelength mode is a \emph{short}-wavelength mode and
is therefore damped quickly by Gauss–Seidel.  Multigrid exploits this
by alternating between fine and coarse grids.

%%──────────────────────────────────────────────────────────────────
\subsection{The Multigrid Algorithm}\label{subsec:lec08:multigrid}

\subsubsection*{Physical Motivation}
We want to solve the linear problem
\begin{equation}
  \mathcal{L}_h u_h = f_h,
  \label{eq:lec08:fine-problem}
\end{equation}
where $\mathcal{L}_h$ is the discrete elliptic operator on a grid of
spacing $h$ (for example $\mathcal{L}_h = \Delta^2_h$ for the Poisson
equation, $f_h = \rho_h$).

\paragraph{Key quantities.}
Let $\tilde{u}_h$ denote the current approximation.  Define:
\begin{align}
  v_h &= u_h - \tilde{u}_h \quad \text{(error)},
  \label{eq:lec08:error}\\
  d_h &= \mathcal{L}_h\tilde{u}_h - f_h \quad \text{(defect / residual)}.
  \label{eq:lec08:defect}
\end{align}
Because $\mathcal{L}_h$ is linear,
\begin{equation}
  \mathcal{L}_h v_h = -d_h.
  \label{eq:lec08:error-equation}
\end{equation}

\textit{Physical significance:} instead of solving for the full
solution $u_h$, we need only solve for the error correction $v_h$,
which satisfies the same type of equation driven by the current
residual.

\paragraph{Coarse-grid correction.}
Pass to a grid of spacing $H=2h$.  Define the \emph{restriction}
operator $\mathcal{R}$ that maps a fine-grid function to the coarse grid:
\begin{equation}
  d_H = \mathcal{R}\,d_h.
  \label{eq:lec08:restriction}
\end{equation}
Solve the coarse-grid error equation approximately to obtain
$\tilde{v}_H$, then \emph{prolong} the correction back to the fine
grid using the prolongation operator $\mathcal{P}$:
\begin{equation}
  \tilde{v}_h = \mathcal{P}\,\tilde{v}_H.
  \label{eq:lec08:prolongation}
\end{equation}
Update the fine-grid approximation:
\begin{equation}
  \boxed{\tilde{u}_h^{\mathrm{new}} = \tilde{u}_h + \tilde{v}_h}
  \label{eq:lec08:update}
\end{equation}

\textit{Physical significance:} the coarse grid corrects the
smooth (long-wavelength) components of the error at low cost because
on the coarse grid these components are short-wavelength and are
smoothed efficiently.

\paragraph{One multigrid cycle.}
\begin{enumerate}
  \item \textbf{Pre-smoothing}: apply $i_1$ Gauss–Seidel iterations
    to Eq.~\eqref{eq:lec08:fine-problem} to obtain $\tilde{u}_h$.
  \item \textbf{Compute defect}: $d_h = \mathcal{L}_h\tilde{u}_h - f_h$.
  \item \textbf{Restrict}: $d_H = \mathcal{R}\,d_h$.
  \item \textbf{Solve coarse}: obtain $\tilde{v}_H\approx -\mathcal{L}_H^{-1}d_H$
    (recursively or by direct solve at the coarsest level, marked
    E in the cycle diagrams).
  \item \textbf{Prolong}: $\tilde{v}_h = \mathcal{P}\,\tilde{v}_H$.
  \item \textbf{Update}: $\tilde{u}_h^{\mathrm{new}}=\tilde{u}_h+\tilde{v}_h$.
  \item \textbf{Post-smoothing}: apply $i_1$ Gauss–Seidel iterations.
\end{enumerate}

The full recursion over $L$ grid levels is a \emph{V-cycle} (one visit
to each coarser level) or a \emph{W-cycle} (two visits).  For an
$N\times N$ grid:
\begin{equation}
  \boxed{\text{Total cost per multigrid cycle} = \bigO{N^2}}
  \label{eq:lec08:multigrid-cost}
\end{equation}

\textit{Physical significance:} because each cycle reduces all error
components by a fixed factor independent of $N$, a fixed number of
cycles suffices for any prescribed accuracy, making the overall
algorithm $\bigO{N^2}$ — optimal up to the unavoidable cost of writing
the answer.

%%──────────────────────────────────────────────────────────────────
\subsection{Classification of Partial Differential Equations}
\label{subsec:lec08:classification}

\subsubsection*{Physical Motivation}
Many fundamental laws of physics — electrostatics, heat conduction,
fluid mechanics, quantum mechanics — are expressed as PDEs.  The
numerical strategy (explicit vs implicit time-stepping, type of
boundary condition, stability criterion) depends strongly on which
class a PDE belongs to.

A second-order PDE in two variables is classified according to the sign
of its discriminant.  The three canonical types and their physical
representatives are:

\begin{itemize}
  \item \textbf{Elliptic} (both eigenvalues same sign): stationary
    problems. Prototype: Poisson equation
    \begin{equation}
      \boxed{\frac{\partial^2 u}{\partial x^2}
        +\frac{\partial^2 u}{\partial y^2}
        = \rho(x,y)}
      \label{eq:lec08:poisson}
    \end{equation}
    \textit{Physical significance:} describes electrostatic potential
    with charge density $\rho$, steady-state temperature, etc.  The
    solution at each point is influenced by the entire domain —
    ``action at a distance''.

  \item \textbf{Parabolic} (one zero eigenvalue): diffusion/heat
    problems with a first-order time derivative. Prototype:
    \begin{equation}
      \boxed{\frac{\partial u}{\partial t}
        = \frac{\partial}{\partial x}
          \!\left(D\frac{\partial u}{\partial x}\right)}
      \label{eq:lec08:diffusion}
    \end{equation}
    \textit{Physical significance:} information propagates at
    infinite speed; the solution smooths out irregularities
    exponentially fast.

  \item \textbf{Hyperbolic} (eigenvalues of opposite sign): wave
    problems with a second-order time derivative. Prototype:
    \begin{equation}
      \boxed{\frac{\partial^2 u}{\partial t^2}
        = v^2\frac{\partial^2 u}{\partial x^2}}
      \label{eq:lec08:wave}
    \end{equation}
    \textit{Physical significance:} information propagates at finite
    speed $v$; wave fronts and shocks can form.
\end{itemize}

%%──────────────────────────────────────────────────────────────────
\subsection{Boundary Conditions}\label{subsec:lec08:bcs}

\subsubsection*{Physical Motivation}
For an elliptic PDE defined in a domain, the solution is uniquely
determined once appropriate boundary conditions are supplied.
Two standard types are:
\begin{itemize}
  \item \textbf{Dirichlet}: the function value is specified on the
    boundary, $u\big|_{\partial\Omega}=g$.
  \item \textbf{Neumann}: the normal derivative is specified,
    $\partial u/\partial n\big|_{\partial\Omega}=g$.
\end{itemize}
For \emph{initial-value} problems (parabolic and hyperbolic) the
function and, where relevant, its time derivative are given on the
initial surface and boundary conditions are applied on the spatial
boundary.

%%──────────────────────────────────────────────────────────────────
\subsection{Discretising the 2D Poisson Equation}
\label{subsec:lec08:poisson-discretisation}

\subsubsection*{Physical Motivation}
We want to solve
$\partial^2_x u + \partial^2_y u = \rho(x,y)$
numerically on a square domain $[0,1]^2$.  The approach is to
replace the partial derivatives by finite differences on a uniform
grid.

\paragraph{Grid.}
Divide each axis into $N$ equal intervals of width $h=1/N$.  Let
\begin{equation}
  u_{i,j} = u(x_i,y_j),
  \qquad
  \rho_{i,j} = \rho(x_i,y_j),
  \qquad
  x_i = ih,\;y_j = jh.
  \label{eq:lec08:grid}
\end{equation}
Applying the symmetric second-derivative formula in both directions
gives the \emph{five-point stencil}:
\begin{equation}
  \boxed{
  u_{i+1,j}+u_{i-1,j}+u_{i,j+1}+u_{i,j-1}-4u_{i,j}
  = h^2\rho_{i,j}}
  \label{eq:lec08:five-point}
\end{equation}
with local truncation error $\bigO{h^2}$.

\textit{Physical significance:} each interior grid point produces one
equation; with $N\times N$ unknowns and Dirichlet boundary conditions
on all sides, the system is square with $N^2$ equations and $N^2$
unknowns.  In matrix form it is a sparse, banded, negative-definite
system.  Even for $N=100$ this is a $10\,000\times10\,000$ matrix —
direct Gaussian elimination costs $\bigO{N^6}$ which is too
expensive; iterative solvers are essential.

%%──────────────────────────────────────────────────────────────────
\subsection{Iterative Solvers for the Poisson System}
\label{subsec:lec08:iterative}

\subsubsection*{Physical Motivation}
The five-point system Eq.~\eqref{eq:lec08:five-point} can be solved
by rearranging for $u_{i,j}$:
\begin{equation}
  u_{i,j}
  = \tfrac{1}{4}\!\left(
      u_{i+1,j}+u_{i-1,j}+u_{i,j+1}+u_{i,j-1}
    \right)
    -\tfrac{h^2}{4}\rho_{i,j}.
  \label{eq:lec08:jacobi-update}
\end{equation}
This is a \emph{local average} — the solution at each point is the
mean of its four neighbours minus a source term.  Iterating this
formula gives three increasingly powerful methods.

\paragraph{Jacobi iteration.}
All grid values are updated simultaneously using only old values:
\begin{equation}
  u_{i,j}^{n+1}
  = \tfrac{1}{4}\!\left(
      u_{i+1,j}^{n}+u_{i-1,j}^{n}+u_{i,j+1}^{n}+u_{i,j-1}^{n}
    \right)
    -\tfrac{h^2}{4}\rho_{i,j}.
  \label{eq:lec08:jacobi}
\end{equation}
Convergence requires $\tfrac{1}{2}pN^2$ iterations for precision
$10^{-p}$.

\paragraph{Gauss–Seidel iteration with red-black ordering.}
Update $u_{i,j}$ using the most recent values of its neighbours.
Using a \emph{red-black} (checkerboard) ordering — first all
$(i+j)$-even nodes (``red''), then all $(i+j)$-odd nodes (``black'')
— the red nodes depend only on black neighbours and vice versa, so each
colour can be updated simultaneously:
\begin{equation}
  u_{i,j}^{n+1}
  = \tfrac{1}{4}\!\left(
      u_{i+1,j}^{n+1}+u_{i-1,j}^{n+1}
      +u_{i,j+1}^{n+1}+u_{i,j-1}^{n+1}
    \right)
    -\tfrac{h^2}{4}\rho_{i,j}
  \label{eq:lec08:gauss-seidel}
\end{equation}
(where newly updated values are used as soon as they are
available). Gauss–Seidel converges roughly twice as fast as Jacobi.

\paragraph{Successive Over-Relaxation (SOR).}
Compute the Gauss–Seidel update $u^*_{i,j}$ and then over-shoot:
\begin{equation}
  u_{i,j}^{n+1}
  = w\,u^*_{i,j} + (1-w)\,u_{i,j}^{n},
  \qquad 1 < w < 2.
  \label{eq:lec08:sor}
\end{equation}
\begin{equation}
  \boxed{w_{\mathrm{opt}}
    \approx \frac{2}{1+\sin(\pi/N)}}
  \label{eq:lec08:sor-optimal}
\end{equation}

\textit{Physical significance:} with optimal $w$, SOR needs only
$\bigO{N}$ iterations instead of $\bigO{N^2}$, reducing the total
cost for the Poisson equation to $\bigO{N^3}$.  For $w=1$ one
recovers plain Gauss–Seidel; $w=2$ is unstable.  Under-relaxation
($0<w<1$) is used when the plain Gauss–Seidel iterations diverge
(e.g.\ for nonlinear problems).

%%──────────────────────────────────────────────────────────────────
\subsection*{Recommended Reading}
\begin{itemize}
  \item \textit{Computational Physics}, Mark Newman --- Ch.\ 9.1--9.3
    (relaxation, Gauss–Seidel, over-relaxation for PDEs).
  \item \textit{Numerical Recipes}, W.H. Press et al.\ --- Ch.\ 19.1
    (relaxation methods for PDEs) and Ch.\ 20.6 (multigrid methods).
  \item K.W. Morton \& D.F. Mayers, \textit{Numerical Solution of
    Partial Differential Equations} --- Ch.\ 2--3 (classification and
    elliptic equations).
\end{itemize}
